\documentclass[11pt]{article}
\usepackage[margin=2.3cm]{geometry}
\usepackage{amsmath,amssymb,amsthm,booktabs,enumitem}
\usepackage[hidelinks]{hyperref}
\title{Peer review: Rick's \emph{Strict reciprocity for Schubert varieties in the Lagrangian
Grassmannian} (FPSAC cut \texttt{6dbecb8} $=$ submission build \texttt{9fa6b0c})}
\author{Clio}
\date{2026-11-11}
\makeatletter
%% MIRROR of proofs/protocol-header.tex (clio-vega/proofs), copied 2026-10-16.
%% reviews/ is a SEPARATE REPO (clio-vega/rick-review) and must build standalone for
%% Rick and Robin, so a symlink or TEXINPUTS dependency would break for them. Hence a
%% copy -- but a VERIFIABLE one, because a silent fork is the same fault in new clothes:
%%   source md5 at copy time: f828569e31971b23bc5726db5d23ecef
%%   drift check:  md5sum ../proofs/protocol-header.tex   # compare with the line above
%% If they differ, re-copy from proofs/ -- that file is the source of truth.
%%
%% protocol-header.tex -- PROTOCOL §2.3 page-1 block, as a TEMPLATE not an instance.
%%
%% PROTOCOL §2.3: "Every PDF must state, on its first page: author, date,
%% recipient, a title, and the commit hash of the work-in-progress push it
%% corresponds to.  A PDF with no commit hash cannot be checked against the
%% record and should be treated as a draft opinion, not a result."
%%
%% Written 2026-10-16 because the fault had recurred twice and I had twice
%% fixed only the instance:
%%   - 2026-10-14: the Q473 PDF went to Rick with no recipient and no hash.
%%   - 2026-10-14: the Q460 PDF went to Rick saying "Recipient: Robin Langer".
%%     He noticed, and took the commit hash as authority instead.
%% Both times I repaired the one file and wrote "fix the template, not just the
%% instance" in my notes.  This file is that fix.  \input it; do not retype it.
%%
%% USAGE -- TWO PLACES, NOT ONE.  The \input goes in the PREAMBLE inside
%% \makeatletter; only \protocolheader goes after \maketitle:
%%
%%   \makeatletter              <- REQUIRED: \protocol@date has an @ in its name
%%   %% MIRROR of proofs/protocol-header.tex (clio-vega/proofs), copied 2026-10-16.
%% reviews/ is a SEPARATE REPO (clio-vega/rick-review) and must build standalone for
%% Rick and Robin, so a symlink or TEXINPUTS dependency would break for them. Hence a
%% copy -- but a VERIFIABLE one, because a silent fork is the same fault in new clothes:
%%   source md5 at copy time: f828569e31971b23bc5726db5d23ecef
%%   drift check:  md5sum ../proofs/protocol-header.tex   # compare with the line above
%% If they differ, re-copy from proofs/ -- that file is the source of truth.
%%
%% protocol-header.tex -- PROTOCOL §2.3 page-1 block, as a TEMPLATE not an instance.
%%
%% PROTOCOL §2.3: "Every PDF must state, on its first page: author, date,
%% recipient, a title, and the commit hash of the work-in-progress push it
%% corresponds to.  A PDF with no commit hash cannot be checked against the
%% record and should be treated as a draft opinion, not a result."
%%
%% Written 2026-10-16 because the fault had recurred twice and I had twice
%% fixed only the instance:
%%   - 2026-10-14: the Q473 PDF went to Rick with no recipient and no hash.
%%   - 2026-10-14: the Q460 PDF went to Rick saying "Recipient: Robin Langer".
%%     He noticed, and took the commit hash as authority instead.
%% Both times I repaired the one file and wrote "fix the template, not just the
%% instance" in my notes.  This file is that fix.  \input it; do not retype it.
%%
%% USAGE -- TWO PLACES, NOT ONE.  The \input goes in the PREAMBLE inside
%% \makeatletter; only \protocolheader goes after \maketitle:
%%
%%   \makeatletter              <- REQUIRED: \protocol@date has an @ in its name
%%   %% MIRROR of proofs/protocol-header.tex (clio-vega/proofs), copied 2026-10-16.
%% reviews/ is a SEPARATE REPO (clio-vega/rick-review) and must build standalone for
%% Rick and Robin, so a symlink or TEXINPUTS dependency would break for them. Hence a
%% copy -- but a VERIFIABLE one, because a silent fork is the same fault in new clothes:
%%   source md5 at copy time: f828569e31971b23bc5726db5d23ecef
%%   drift check:  md5sum ../proofs/protocol-header.tex   # compare with the line above
%% If they differ, re-copy from proofs/ -- that file is the source of truth.
%%
%% protocol-header.tex -- PROTOCOL §2.3 page-1 block, as a TEMPLATE not an instance.
%%
%% PROTOCOL §2.3: "Every PDF must state, on its first page: author, date,
%% recipient, a title, and the commit hash of the work-in-progress push it
%% corresponds to.  A PDF with no commit hash cannot be checked against the
%% record and should be treated as a draft opinion, not a result."
%%
%% Written 2026-10-16 because the fault had recurred twice and I had twice
%% fixed only the instance:
%%   - 2026-10-14: the Q473 PDF went to Rick with no recipient and no hash.
%%   - 2026-10-14: the Q460 PDF went to Rick saying "Recipient: Robin Langer".
%%     He noticed, and took the commit hash as authority instead.
%% Both times I repaired the one file and wrote "fix the template, not just the
%% instance" in my notes.  This file is that fix.  \input it; do not retype it.
%%
%% USAGE -- TWO PLACES, NOT ONE.  The \input goes in the PREAMBLE inside
%% \makeatletter; only \protocolheader goes after \maketitle:
%%
%%   \makeatletter              <- REQUIRED: \protocol@date has an @ in its name
%%   \input{protocol-header}    <- PREAMBLE, while \@date is still live
%%   \makeatother
%%   ...
%%   \begin{document}
%%   \maketitle
%%   \protocolheader{Rick (grandpa-rick)}{clio-vega/proofs}{60e4ce3}
%%
%% CORRECTED 2026-10-16 (peer review): this block previously read "USAGE,
%% immediately after \maketitle: \input{protocol-header} \protocolheader{...}",
%% which is WRONG TWICE -- it omits \makeatletter, and it puts the \input after
%% \maketitle, by which point \@date has been cleared.  I followed it verbatim
%% while building the PROVE-241 review and page 1 printed a literal "@@date"
%% plus a stray "datedate", with pdflatex exiting 0.  The three proofs/*.tex that
%% use this file were all written with the CORRECT pattern (\makeatletter +
%% \input in the preamble), so their PDFs are fine and were verified by
%% pdftotext -- the defect was in this instruction, not in the macro.
%% An instruction that encodes a fault renews it in every session that obeys it.
%%
%% The recipient argument is MANDATORY and has no default on purpose: a default
%% is what produced "Recipient: Robin Langer" on a PDF addressed to Rick.

%% \maketitle CLEARS \@date (and \@author, \@title) in the article class, so a
%% macro called AFTER \maketitle that reads \@date prints an EMPTY row.  Caught
%% 2026-10-16 only because I ran pdftotext on page 1 instead of trusting rc=0:
%% the build was clean and the Date field was blank.  Hence: snapshot the date
%% HERE, at \input time in the preamble, where \@date is still live.
\let\protocol@date\@date

\newcommand{\protocolheader}[3]{%
  \begin{center}
  \fbox{\begin{minipage}{0.92\textwidth}
  \small
  \textbf{PROTOCOL §2.3 record.}\\[1mm]
  \begin{tabular}{@{}ll@{}}
  Author     & Clio \\
  Date       & \protocol@date \\
  Recipient  & #1 \\
  Repository & \texttt{#2} \\
  Commit     & \texttt{#3} \\
  \end{tabular}\\[1mm]
  This PDF corresponds to the pushed commit above.  Where this text and the
  commit disagree, \textbf{the commit is authoritative}.
  \end{minipage}}
  \end{center}
  \vspace{2mm}%
}
    <- PREAMBLE, while \@date is still live
%%   \makeatother
%%   ...
%%   \begin{document}
%%   \maketitle
%%   \protocolheader{Rick (grandpa-rick)}{clio-vega/proofs}{60e4ce3}
%%
%% CORRECTED 2026-10-16 (peer review): this block previously read "USAGE,
%% immediately after \maketitle: %% MIRROR of proofs/protocol-header.tex (clio-vega/proofs), copied 2026-10-16.
%% reviews/ is a SEPARATE REPO (clio-vega/rick-review) and must build standalone for
%% Rick and Robin, so a symlink or TEXINPUTS dependency would break for them. Hence a
%% copy -- but a VERIFIABLE one, because a silent fork is the same fault in new clothes:
%%   source md5 at copy time: f828569e31971b23bc5726db5d23ecef
%%   drift check:  md5sum ../proofs/protocol-header.tex   # compare with the line above
%% If they differ, re-copy from proofs/ -- that file is the source of truth.
%%
%% protocol-header.tex -- PROTOCOL §2.3 page-1 block, as a TEMPLATE not an instance.
%%
%% PROTOCOL §2.3: "Every PDF must state, on its first page: author, date,
%% recipient, a title, and the commit hash of the work-in-progress push it
%% corresponds to.  A PDF with no commit hash cannot be checked against the
%% record and should be treated as a draft opinion, not a result."
%%
%% Written 2026-10-16 because the fault had recurred twice and I had twice
%% fixed only the instance:
%%   - 2026-10-14: the Q473 PDF went to Rick with no recipient and no hash.
%%   - 2026-10-14: the Q460 PDF went to Rick saying "Recipient: Robin Langer".
%%     He noticed, and took the commit hash as authority instead.
%% Both times I repaired the one file and wrote "fix the template, not just the
%% instance" in my notes.  This file is that fix.  \input it; do not retype it.
%%
%% USAGE -- TWO PLACES, NOT ONE.  The \input goes in the PREAMBLE inside
%% \makeatletter; only \protocolheader goes after \maketitle:
%%
%%   \makeatletter              <- REQUIRED: \protocol@date has an @ in its name
%%   \input{protocol-header}    <- PREAMBLE, while \@date is still live
%%   \makeatother
%%   ...
%%   \begin{document}
%%   \maketitle
%%   \protocolheader{Rick (grandpa-rick)}{clio-vega/proofs}{60e4ce3}
%%
%% CORRECTED 2026-10-16 (peer review): this block previously read "USAGE,
%% immediately after \maketitle: \input{protocol-header} \protocolheader{...}",
%% which is WRONG TWICE -- it omits \makeatletter, and it puts the \input after
%% \maketitle, by which point \@date has been cleared.  I followed it verbatim
%% while building the PROVE-241 review and page 1 printed a literal "@@date"
%% plus a stray "datedate", with pdflatex exiting 0.  The three proofs/*.tex that
%% use this file were all written with the CORRECT pattern (\makeatletter +
%% \input in the preamble), so their PDFs are fine and were verified by
%% pdftotext -- the defect was in this instruction, not in the macro.
%% An instruction that encodes a fault renews it in every session that obeys it.
%%
%% The recipient argument is MANDATORY and has no default on purpose: a default
%% is what produced "Recipient: Robin Langer" on a PDF addressed to Rick.

%% \maketitle CLEARS \@date (and \@author, \@title) in the article class, so a
%% macro called AFTER \maketitle that reads \@date prints an EMPTY row.  Caught
%% 2026-10-16 only because I ran pdftotext on page 1 instead of trusting rc=0:
%% the build was clean and the Date field was blank.  Hence: snapshot the date
%% HERE, at \input time in the preamble, where \@date is still live.
\let\protocol@date\@date

\newcommand{\protocolheader}[3]{%
  \begin{center}
  \fbox{\begin{minipage}{0.92\textwidth}
  \small
  \textbf{PROTOCOL §2.3 record.}\\[1mm]
  \begin{tabular}{@{}ll@{}}
  Author     & Clio \\
  Date       & \protocol@date \\
  Recipient  & #1 \\
  Repository & \texttt{#2} \\
  Commit     & \texttt{#3} \\
  \end{tabular}\\[1mm]
  This PDF corresponds to the pushed commit above.  Where this text and the
  commit disagree, \textbf{the commit is authoritative}.
  \end{minipage}}
  \end{center}
  \vspace{2mm}%
}
 \protocolheader{...}",
%% which is WRONG TWICE -- it omits \makeatletter, and it puts the \input after
%% \maketitle, by which point \@date has been cleared.  I followed it verbatim
%% while building the PROVE-241 review and page 1 printed a literal "@@date"
%% plus a stray "datedate", with pdflatex exiting 0.  The three proofs/*.tex that
%% use this file were all written with the CORRECT pattern (\makeatletter +
%% \input in the preamble), so their PDFs are fine and were verified by
%% pdftotext -- the defect was in this instruction, not in the macro.
%% An instruction that encodes a fault renews it in every session that obeys it.
%%
%% The recipient argument is MANDATORY and has no default on purpose: a default
%% is what produced "Recipient: Robin Langer" on a PDF addressed to Rick.

%% \maketitle CLEARS \@date (and \@author, \@title) in the article class, so a
%% macro called AFTER \maketitle that reads \@date prints an EMPTY row.  Caught
%% 2026-10-16 only because I ran pdftotext on page 1 instead of trusting rc=0:
%% the build was clean and the Date field was blank.  Hence: snapshot the date
%% HERE, at \input time in the preamble, where \@date is still live.
\let\protocol@date\@date

\newcommand{\protocolheader}[3]{%
  \begin{center}
  \fbox{\begin{minipage}{0.92\textwidth}
  \small
  \textbf{PROTOCOL §2.3 record.}\\[1mm]
  \begin{tabular}{@{}ll@{}}
  Author     & Clio \\
  Date       & \protocol@date \\
  Recipient  & #1 \\
  Repository & \texttt{#2} \\
  Commit     & \texttt{#3} \\
  \end{tabular}\\[1mm]
  This PDF corresponds to the pushed commit above.  Where this text and the
  commit disagree, \textbf{the commit is authoritative}.
  \end{minipage}}
  \end{center}
  \vspace{2mm}%
}
    <- PREAMBLE, while \@date is still live
%%   \makeatother
%%   ...
%%   \begin{document}
%%   \maketitle
%%   \protocolheader{Rick (grandpa-rick)}{clio-vega/proofs}{60e4ce3}
%%
%% CORRECTED 2026-10-16 (peer review): this block previously read "USAGE,
%% immediately after \maketitle: %% MIRROR of proofs/protocol-header.tex (clio-vega/proofs), copied 2026-10-16.
%% reviews/ is a SEPARATE REPO (clio-vega/rick-review) and must build standalone for
%% Rick and Robin, so a symlink or TEXINPUTS dependency would break for them. Hence a
%% copy -- but a VERIFIABLE one, because a silent fork is the same fault in new clothes:
%%   source md5 at copy time: f828569e31971b23bc5726db5d23ecef
%%   drift check:  md5sum ../proofs/protocol-header.tex   # compare with the line above
%% If they differ, re-copy from proofs/ -- that file is the source of truth.
%%
%% protocol-header.tex -- PROTOCOL §2.3 page-1 block, as a TEMPLATE not an instance.
%%
%% PROTOCOL §2.3: "Every PDF must state, on its first page: author, date,
%% recipient, a title, and the commit hash of the work-in-progress push it
%% corresponds to.  A PDF with no commit hash cannot be checked against the
%% record and should be treated as a draft opinion, not a result."
%%
%% Written 2026-10-16 because the fault had recurred twice and I had twice
%% fixed only the instance:
%%   - 2026-10-14: the Q473 PDF went to Rick with no recipient and no hash.
%%   - 2026-10-14: the Q460 PDF went to Rick saying "Recipient: Robin Langer".
%%     He noticed, and took the commit hash as authority instead.
%% Both times I repaired the one file and wrote "fix the template, not just the
%% instance" in my notes.  This file is that fix.  \input it; do not retype it.
%%
%% USAGE -- TWO PLACES, NOT ONE.  The \input goes in the PREAMBLE inside
%% \makeatletter; only \protocolheader goes after \maketitle:
%%
%%   \makeatletter              <- REQUIRED: \protocol@date has an @ in its name
%%   %% MIRROR of proofs/protocol-header.tex (clio-vega/proofs), copied 2026-10-16.
%% reviews/ is a SEPARATE REPO (clio-vega/rick-review) and must build standalone for
%% Rick and Robin, so a symlink or TEXINPUTS dependency would break for them. Hence a
%% copy -- but a VERIFIABLE one, because a silent fork is the same fault in new clothes:
%%   source md5 at copy time: f828569e31971b23bc5726db5d23ecef
%%   drift check:  md5sum ../proofs/protocol-header.tex   # compare with the line above
%% If they differ, re-copy from proofs/ -- that file is the source of truth.
%%
%% protocol-header.tex -- PROTOCOL §2.3 page-1 block, as a TEMPLATE not an instance.
%%
%% PROTOCOL §2.3: "Every PDF must state, on its first page: author, date,
%% recipient, a title, and the commit hash of the work-in-progress push it
%% corresponds to.  A PDF with no commit hash cannot be checked against the
%% record and should be treated as a draft opinion, not a result."
%%
%% Written 2026-10-16 because the fault had recurred twice and I had twice
%% fixed only the instance:
%%   - 2026-10-14: the Q473 PDF went to Rick with no recipient and no hash.
%%   - 2026-10-14: the Q460 PDF went to Rick saying "Recipient: Robin Langer".
%%     He noticed, and took the commit hash as authority instead.
%% Both times I repaired the one file and wrote "fix the template, not just the
%% instance" in my notes.  This file is that fix.  \input it; do not retype it.
%%
%% USAGE -- TWO PLACES, NOT ONE.  The \input goes in the PREAMBLE inside
%% \makeatletter; only \protocolheader goes after \maketitle:
%%
%%   \makeatletter              <- REQUIRED: \protocol@date has an @ in its name
%%   \input{protocol-header}    <- PREAMBLE, while \@date is still live
%%   \makeatother
%%   ...
%%   \begin{document}
%%   \maketitle
%%   \protocolheader{Rick (grandpa-rick)}{clio-vega/proofs}{60e4ce3}
%%
%% CORRECTED 2026-10-16 (peer review): this block previously read "USAGE,
%% immediately after \maketitle: \input{protocol-header} \protocolheader{...}",
%% which is WRONG TWICE -- it omits \makeatletter, and it puts the \input after
%% \maketitle, by which point \@date has been cleared.  I followed it verbatim
%% while building the PROVE-241 review and page 1 printed a literal "@@date"
%% plus a stray "datedate", with pdflatex exiting 0.  The three proofs/*.tex that
%% use this file were all written with the CORRECT pattern (\makeatletter +
%% \input in the preamble), so their PDFs are fine and were verified by
%% pdftotext -- the defect was in this instruction, not in the macro.
%% An instruction that encodes a fault renews it in every session that obeys it.
%%
%% The recipient argument is MANDATORY and has no default on purpose: a default
%% is what produced "Recipient: Robin Langer" on a PDF addressed to Rick.

%% \maketitle CLEARS \@date (and \@author, \@title) in the article class, so a
%% macro called AFTER \maketitle that reads \@date prints an EMPTY row.  Caught
%% 2026-10-16 only because I ran pdftotext on page 1 instead of trusting rc=0:
%% the build was clean and the Date field was blank.  Hence: snapshot the date
%% HERE, at \input time in the preamble, where \@date is still live.
\let\protocol@date\@date

\newcommand{\protocolheader}[3]{%
  \begin{center}
  \fbox{\begin{minipage}{0.92\textwidth}
  \small
  \textbf{PROTOCOL §2.3 record.}\\[1mm]
  \begin{tabular}{@{}ll@{}}
  Author     & Clio \\
  Date       & \protocol@date \\
  Recipient  & #1 \\
  Repository & \texttt{#2} \\
  Commit     & \texttt{#3} \\
  \end{tabular}\\[1mm]
  This PDF corresponds to the pushed commit above.  Where this text and the
  commit disagree, \textbf{the commit is authoritative}.
  \end{minipage}}
  \end{center}
  \vspace{2mm}%
}
    <- PREAMBLE, while \@date is still live
%%   \makeatother
%%   ...
%%   \begin{document}
%%   \maketitle
%%   \protocolheader{Rick (grandpa-rick)}{clio-vega/proofs}{60e4ce3}
%%
%% CORRECTED 2026-10-16 (peer review): this block previously read "USAGE,
%% immediately after \maketitle: %% MIRROR of proofs/protocol-header.tex (clio-vega/proofs), copied 2026-10-16.
%% reviews/ is a SEPARATE REPO (clio-vega/rick-review) and must build standalone for
%% Rick and Robin, so a symlink or TEXINPUTS dependency would break for them. Hence a
%% copy -- but a VERIFIABLE one, because a silent fork is the same fault in new clothes:
%%   source md5 at copy time: f828569e31971b23bc5726db5d23ecef
%%   drift check:  md5sum ../proofs/protocol-header.tex   # compare with the line above
%% If they differ, re-copy from proofs/ -- that file is the source of truth.
%%
%% protocol-header.tex -- PROTOCOL §2.3 page-1 block, as a TEMPLATE not an instance.
%%
%% PROTOCOL §2.3: "Every PDF must state, on its first page: author, date,
%% recipient, a title, and the commit hash of the work-in-progress push it
%% corresponds to.  A PDF with no commit hash cannot be checked against the
%% record and should be treated as a draft opinion, not a result."
%%
%% Written 2026-10-16 because the fault had recurred twice and I had twice
%% fixed only the instance:
%%   - 2026-10-14: the Q473 PDF went to Rick with no recipient and no hash.
%%   - 2026-10-14: the Q460 PDF went to Rick saying "Recipient: Robin Langer".
%%     He noticed, and took the commit hash as authority instead.
%% Both times I repaired the one file and wrote "fix the template, not just the
%% instance" in my notes.  This file is that fix.  \input it; do not retype it.
%%
%% USAGE -- TWO PLACES, NOT ONE.  The \input goes in the PREAMBLE inside
%% \makeatletter; only \protocolheader goes after \maketitle:
%%
%%   \makeatletter              <- REQUIRED: \protocol@date has an @ in its name
%%   \input{protocol-header}    <- PREAMBLE, while \@date is still live
%%   \makeatother
%%   ...
%%   \begin{document}
%%   \maketitle
%%   \protocolheader{Rick (grandpa-rick)}{clio-vega/proofs}{60e4ce3}
%%
%% CORRECTED 2026-10-16 (peer review): this block previously read "USAGE,
%% immediately after \maketitle: \input{protocol-header} \protocolheader{...}",
%% which is WRONG TWICE -- it omits \makeatletter, and it puts the \input after
%% \maketitle, by which point \@date has been cleared.  I followed it verbatim
%% while building the PROVE-241 review and page 1 printed a literal "@@date"
%% plus a stray "datedate", with pdflatex exiting 0.  The three proofs/*.tex that
%% use this file were all written with the CORRECT pattern (\makeatletter +
%% \input in the preamble), so their PDFs are fine and were verified by
%% pdftotext -- the defect was in this instruction, not in the macro.
%% An instruction that encodes a fault renews it in every session that obeys it.
%%
%% The recipient argument is MANDATORY and has no default on purpose: a default
%% is what produced "Recipient: Robin Langer" on a PDF addressed to Rick.

%% \maketitle CLEARS \@date (and \@author, \@title) in the article class, so a
%% macro called AFTER \maketitle that reads \@date prints an EMPTY row.  Caught
%% 2026-10-16 only because I ran pdftotext on page 1 instead of trusting rc=0:
%% the build was clean and the Date field was blank.  Hence: snapshot the date
%% HERE, at \input time in the preamble, where \@date is still live.
\let\protocol@date\@date

\newcommand{\protocolheader}[3]{%
  \begin{center}
  \fbox{\begin{minipage}{0.92\textwidth}
  \small
  \textbf{PROTOCOL §2.3 record.}\\[1mm]
  \begin{tabular}{@{}ll@{}}
  Author     & Clio \\
  Date       & \protocol@date \\
  Recipient  & #1 \\
  Repository & \texttt{#2} \\
  Commit     & \texttt{#3} \\
  \end{tabular}\\[1mm]
  This PDF corresponds to the pushed commit above.  Where this text and the
  commit disagree, \textbf{the commit is authoritative}.
  \end{minipage}}
  \end{center}
  \vspace{2mm}%
}
 \protocolheader{...}",
%% which is WRONG TWICE -- it omits \makeatletter, and it puts the \input after
%% \maketitle, by which point \@date has been cleared.  I followed it verbatim
%% while building the PROVE-241 review and page 1 printed a literal "@@date"
%% plus a stray "datedate", with pdflatex exiting 0.  The three proofs/*.tex that
%% use this file were all written with the CORRECT pattern (\makeatletter +
%% \input in the preamble), so their PDFs are fine and were verified by
%% pdftotext -- the defect was in this instruction, not in the macro.
%% An instruction that encodes a fault renews it in every session that obeys it.
%%
%% The recipient argument is MANDATORY and has no default on purpose: a default
%% is what produced "Recipient: Robin Langer" on a PDF addressed to Rick.

%% \maketitle CLEARS \@date (and \@author, \@title) in the article class, so a
%% macro called AFTER \maketitle that reads \@date prints an EMPTY row.  Caught
%% 2026-10-16 only because I ran pdftotext on page 1 instead of trusting rc=0:
%% the build was clean and the Date field was blank.  Hence: snapshot the date
%% HERE, at \input time in the preamble, where \@date is still live.
\let\protocol@date\@date

\newcommand{\protocolheader}[3]{%
  \begin{center}
  \fbox{\begin{minipage}{0.92\textwidth}
  \small
  \textbf{PROTOCOL §2.3 record.}\\[1mm]
  \begin{tabular}{@{}ll@{}}
  Author     & Clio \\
  Date       & \protocol@date \\
  Recipient  & #1 \\
  Repository & \texttt{#2} \\
  Commit     & \texttt{#3} \\
  \end{tabular}\\[1mm]
  This PDF corresponds to the pushed commit above.  Where this text and the
  commit disagree, \textbf{the commit is authoritative}.
  \end{minipage}}
  \end{center}
  \vspace{2mm}%
}
 \protocolheader{...}",
%% which is WRONG TWICE -- it omits \makeatletter, and it puts the \input after
%% \maketitle, by which point \@date has been cleared.  I followed it verbatim
%% while building the PROVE-241 review and page 1 printed a literal "@@date"
%% plus a stray "datedate", with pdflatex exiting 0.  The three proofs/*.tex that
%% use this file were all written with the CORRECT pattern (\makeatletter +
%% \input in the preamble), so their PDFs are fine and were verified by
%% pdftotext -- the defect was in this instruction, not in the macro.
%% An instruction that encodes a fault renews it in every session that obeys it.
%%
%% The recipient argument is MANDATORY and has no default on purpose: a default
%% is what produced "Recipient: Robin Langer" on a PDF addressed to Rick.

%% \maketitle CLEARS \@date (and \@author, \@title) in the article class, so a
%% macro called AFTER \maketitle that reads \@date prints an EMPTY row.  Caught
%% 2026-10-16 only because I ran pdftotext on page 1 instead of trusting rc=0:
%% the build was clean and the Date field was blank.  Hence: snapshot the date
%% HERE, at \input time in the preamble, where \@date is still live.
\let\protocol@date\@date

\newcommand{\protocolheader}[3]{%
  \begin{center}
  \fbox{\begin{minipage}{0.92\textwidth}
  \small
  \textbf{PROTOCOL §2.3 record.}\\[1mm]
  \begin{tabular}{@{}ll@{}}
  Author     & Clio \\
  Date       & \protocol@date \\
  Recipient  & #1 \\
  Repository & \texttt{#2} \\
  Commit     & \texttt{#3} \\
  \end{tabular}\\[1mm]
  This PDF corresponds to the pushed commit above.  Where this text and the
  commit disagree, \textbf{the commit is authoritative}.
  \end{minipage}}
  \end{center}
  \vspace{2mm}%
}

\makeatother
\newcommand{\LG}{\mathrm{LG}}
\newcommand{\cO}{\mathcal{O}}
\begin{document}
\maketitle
\protocolheader{Rick (grandpa-rick), cc Robin}{clio-vega/rick-review}{138efc9dbd}

\paragraph{Target.} \texttt{grandpa-rick/work-in-progress}, \texttt{papers/lg-og-fpsac-cut/main.tex}
at \texttt{6dbecb8} (requested UID 866, 10-28), and the current build
\texttt{papers/lg-og-fpsac-submit/main.tex} at \texttt{9fa6b0c} (UID 902, 11-11). Both hashes
were resolved upstream with \texttt{gh api}; a bogus SHA returned 422 (negative control).
\textbf{The two files are mathematically identical.} I diffed them myself, not on the strength of
the covering note: there are five hunks, covering the header comment, the author/address/date
block, one keyword, the Dehy footnote attribution, and the deletion of the acknowledgements
paragraph and the \verb|\robin{}| placeholder in the AI declaration. No displayed formula, statement
or proof changes. Theorem numbers below are those of the built PDF.

\paragraph{Context.} FPSAC was withdrawn on 11-11 (UID 902), so the 11-13 and 11-15 deadlines
no longer apply. Rick says the review still decides when the paper becomes a publishable result.
Whether and where it goes next is Robin's decision.

\section*{Verdict}

\begin{center}\small
\begin{tabular}{@{}p{5.2cm}p{4.3cm}p{6.3cm}@{}}
\toprule
item & grade & basis\\\midrule
\S3: Lemmas 3.1, 3.2, Prop.~3.3, \textbf{Thm 1.1} (Main Theorem) & \textbf{endorsed}, proved given (G1),(LS) & every step read and checked by hand; independent Demazure engine, $m\le4$: 180/180 (a), 180/180 (b)\\
Lemmas 4.1, 4.2, \textbf{Prop.~4.3} (recursion $(\star)$) & \textbf{endorsed}, proved given (G1),(LS) & by hand; $(\star)$ checked as an identity at 20 points per $\kappa$, including negative and half-integer $x$\\
\textbf{Thm 1.5} (HZ) incl.\ Prop 6.1, Thm 6.2, Lemma 6.3, Cor 6.4 & \textbf{endorsed with corrections} (N1--N3, wording only) & proof line-read; fibre/shift argument re-derived (\S3); every clause checked\\
\textbf{Cor 1.6} (OC) & \textbf{endorsed with correction} N3 & by hand; checked $a=m,\dots,m+4$\\
Thm 1.4 (LG), Rem 5.1, Rem 6.5 & endorsed (now on the $(\star)$ route) & combinatorial claims checked: 64/64, 26/26, 22/22\\
Rem 4.4 (Pieri values) & \textbf{right shape}; given (BR) & $(-1)^{|\theta|-1}V(\theta)$ equals BR Thm 5.5 on $R(\theta)$, 144/144; wrong orientation fails 24/144\\
Rem 6.6 (half-step $T$) & \textbf{computed only} (N4) & no proof printed; weight list incomplete; $s_{x+1/2}=Ts_x$ 224/224\\
AI declaration, ``Review'' item & \textbf{correction needed} (D1--D3) & compared against my sent mail and review files (\S4)\\
\bottomrule
\end{tabular}
\end{center}

\noindent\textbf{No MUST-level mathematical error.} I found nothing that would change a statement.
The corrections are one inaccurate sentence inside a proof (N1), two notational or edge-case slips
(N2, N3), one remark that is stated but not proved (N4), and three places where the AI declaration
describes my earlier reviews less precisely than I wrote them (D1--D3).

\section{\S3, the LS-path proof of Theorem 1.1}

I read every displayed step. I could reconstruct each inference.

\begin{itemize}[leftmargin=*]
\item \textbf{Lemma 3.1.} (a): for $\mu_i=-1$, the positive roots with negative pairing are
$2\varepsilon_i$, $\varepsilon_i-\varepsilon_k$ with $\mu_k=+1$, and $\varepsilon_i+\varepsilon_k$
with $\mu_k=-1$ ($k>i$). That gives $1+(m-i)=m-i+1$~$\checkmark$. (b,c): I redid the three root families. A
cover is exactly a length-one step. $s_{\varepsilon_i-\varepsilon_k}$ replaces part $q$ by $p$, and
it is a cover iff $p=q+1$, which extends a row by an off-diagonal box with pairing $\pm2$. Next,
$s_{\varepsilon_i+\varepsilon_k}$ adds $p+q\ge3$ boxes, so it is never a cover. Last,
$s_{2\varepsilon_i}$ is a cover iff $p=1$: a new diagonal box, with coroot $\varepsilon_i$ and
pairing $\pm1$~$\checkmark$. My engine independently confirms that reduced-word length equals $|\kappa|$
(BFS over sign vectors).
\item \textbf{Lemma 3.2.} Bonds lie in $\{1,2\}$ and scale by $j$, so $2ja\in\mathbb Z$ is
necessary. If $2ja$ is odd, only bond-2 covers qualify, and every saturated chain adds every box
of the difference~$\checkmark$.
\item \textbf{Prop.~3.3.} $g$ is order-reversing because the $D(\tau_i)$ are nested ideals. It is
even on the diagonal because a diagonal box in $D(\tau_i)\setminus D(\tau_{i+1})$ forces $2ja_i$
even. Its inverse is $\tau_i=\{g\ge v_i\}$, $\tau_{r+1}=g^{-1}(2j)$, which may be empty
(direction $e$, allowed). The two maps are mutually inverse~$\checkmark$.
\item \textbf{Proof of Thm 1.1.} After the off-diagonal $\times2$ rescaling, the vertices of
$\mathcal O(P_\kappa)$ have coordinates in $\{0,1,2\}$, so Ehrhart--Macdonald applies to an honest
lattice polytope. The interior description is the strict count $N_\kappa$~$\checkmark$.
\item \textbf{Lemmas 2.1--2.3}, used here. Lemma 2.2: $D_i(\mathfrak m^k)\subseteq\mathfrak m^{k-1}$
via $D_if=(f-s_if)/(1-e^{-\alpha_i})+s_if$~$\checkmark$. Lemma 2.3: the only right descent of a minimal
representative is $s_m$, and $\langle-\omega,\alpha_m^\vee\rangle=-1$~$\checkmark$. Lemma 2.1(a): the
staircase chain to a box of maximal $a+b$ has $r$ boxes and starts at the even-valued $(0,0)$~$\checkmark$.
\end{itemize}
\emph{Inputs.} The main line uses exactly (G1) and (LS). It does \emph{not} use (G0) or (BR); the
Chevalley formula appears only in Rem 4.4. This is a real improvement on Theorem I/Z, which I
endorsed in October only given (G0),(G1),(BR). The (LS) locator \cite[\S5.2]{Lit94} is declared as
quoted second-hand. I have not read Invent.\ Math.\ 116 either, so that locator is unverified on my
side too.

\section{\S4: Lemmas 4.1, 4.2 and Proposition 4.3}

\textbf{Lemma 4.1.} The map is forced on non-maximal and non-minimal elements, and only isolated
non-$D$ elements are free. For the last claim, $(a,b)<(c,c)$ with $a\le b$ forces $a\le c-1$,
hence $(a,b)\le(c-1,c)\le(c,c)$, and $(c-1,c)\in\theta$ because $\theta$ is an up-set of an ideal
$\checkmark$. The closed form matches brute force on all 144 pairs at $m\le4$.

\textbf{Lemma 4.2.} $\theta=f^{-1}\{2j,2j+1\}$ is an up-set, and $2j$ is the only even top value.
Gluing works because nothing in $\theta$ lies below $D(\rho)$~$\checkmark$.

\textbf{Prop.~4.3.} This is Lemma 4.2 together with Thm 1.1(b) for \emph{every}
$\rho\subsetneq\kappa$, including $N_\varnothing=1$, followed by polynomial extension in $x$~$\checkmark$.
Since $(\star)$ is a polynomial identity, I tested it at
$x\in\{-3,-\tfrac52,\dots,\tfrac{13}2\}$: 600/600 at $m\le4$. A lesion replacing $V$ by the \emph{BR-oriented} count on $\theta$ ($1'<1$ read
directly, diagonal $=$ larger letter) fails on 16/30 shapes. So the test does see orientation.

\textbf{Rem 4.4, on the right shape.} My 10-24 F4 said the closed form must be read on $R(\theta)$.
The cut does this correctly. It defines $V$ with the order \emph{reversed} ($e<o$, diagonal $=e$)
directly on $\theta$, which is the transport of BR's rule along the order-reversing $R$. I compared
$(-1)^{|\theta|-1}V(\theta)$ (brute force, the cut's definition) with BR's Thm 5.5 recursion
$\mathcal C(R(\theta),1)$, implemented on 10-24 from the statement in arXiv:1005.2605v1. Result:
144/144 at $m\le4$. The wrong orientation fails 24/144. The worked example for $(3,2)$ matches, and
so does the count of five $\rho$ for $(3,1)$.

\emph{BR locators} (re-verified first-hand today, by reconstructing the shared \texttt{thm} counter
of \texttt{kpieri.tex} v1):
4.4 $=$ the corollary $c^\nu_{\lambda,p}=\mathcal B(\nu/\lambda,p)$; 4.5 $=$
\texttt{lem:vanishB}, whose (b),(c) give support $=$ rook strips with value $(-1)^{|\theta|-1}$;
5.4, 5.5 $=$ \texttt{thm:pieriCrec}; 5.6 $=$ \texttt{cor:pieriCtab}. All land. Rem 4.5 identifies
``rook strip'' with ``antichain of $P_\kappa$''. That holds for a skew shape: $(a,b),(a',b')$ with
$a<a'$, $b<b'$ forces $(a,b')\in\theta$. The bib entry correctly says ``numbering cited from arXiv
v1''.

\section{\S6: Theorem 1.5 (HZ) and Corollary 1.6 (OC)}

\textbf{Prop.~6.1 (sketch).} I redid the two intertwinings. For $i=m$:
$s_m e^{-\omega/2}=e^{-\omega/2}y_m^{-1}$, so
$D_m(e^{-\omega/2}g)=e^{-\omega/2}(1+y_m)^{-1}D^B_mg$. Next,
$s_m(1+y_m)^{-1}=y_m(1+y_m)^{-1}$, which gives $D'_m(\Phi g)=\Phi D_mg$~$\checkmark$. I also checked the
identity directly, as an independent route to the half-integer values ($D_w(F_a)$ at $y=1$
against the interpolated polynomial): 120/120 at $m\le4$, $a\le3$.

\textbf{Thm 6.2.} The word $c(\kappa_\ell)\cdots c(\kappa_1)$ is reduced. $D_q$ is linear over the
$s_q$-invariant $\prod_{i<q}(1+y_i)$ and kills $y_q$. The remaining letters $s_j$, $j\ge p+1$, fix
$\prod_{i<p}(1+y_i)$~$\checkmark$.

\textbf{Lemma 6.3.} (a),(b): peeling the top value gives $o,e,o,\dots,o$, hence
$(A_oA_e)^{a-1}A_o$~$\checkmark$ (checked 174/174). (c): the toggle $\theta\mapsto\theta\triangle\{z\}$ with
$z$ the last box of the bottom row is a well-defined involution on antichain up-sets. Here is why.
If $z\notin\theta$, no $y<z$ lies in $\theta$ (up-set), so $\theta\cup\{z\}$ is an antichain up-set.
It flips $|\nu|$ and keeps $\nu_1$ because $\ell\ge2$~$\checkmark$ (30/30).

\textbf{Proof of Thm 1.5.} I re-derived the shift: $G_L(M;u)=G_L(M-1;u-1)+G_{L+1}(M-1;u-1)$ when
$u_L\ge2$ and $M\ge1$, and the telescoping of $(-\tfrac12)^{L'-L}$ gives
$\Psi_L=G_L+\tfrac12\Psi_{L+1}$~$\checkmark$. Branch $2a\le n$: after $M$ steps every $G$ needs more values
than are available, and $\Psi_{M+1}(0)=(-\tfrac12)^{n-1-M}$, so
$\Phi=\tfrac12\cdot2^{-M}(-\tfrac12)^{n-1-M}=(-1)^n2^{-n}$ ($M$ odd)~$\checkmark$. Branch $2a\ge n$: $n-1$
steps with all coefficients positive end at $\Psi_n=1$~$\checkmark$. On $2a=n$ the two branches agree because
$n$ is even~$\checkmark$. Signs, the bound, integrality and the vanishing criterion all follow as written.

\begin{itemize}[leftmargin=*]
\item[\textbf{N1}] \emph{``Since $g(1,1)$ is even, $u_\beta\ge\beta+1$.''} This is false for
absent boxes $(1,\beta)$ with $\beta>M$, where $u_\beta=M+1\le\beta$. Example: $\kappa=(5)$,
$a=1$ gives $u_\beta=2$ for $\beta=2,3,4$. The proof survives. What the iteration uses at step
$t$ is $u_\beta-t\ge2$ for $\beta\ge t+1$, with $t\le M-1$ (resp.\ $t\le n-2\le M-1$). This holds
because present boxes have $u_\beta\ge\beta+1$ and absent ones have $u_\beta=M+1$. Suggested
wording: ``$u_\beta\ge\min(\beta,M)+1$''. I worked $\kappa=(5)$, $a=1$ through by hand:
$\Phi=-2^{-5}$, as claimed.
\item[\textbf{N2}] $M$ means $2a-1$ in the proof of Thm 1.5, $2a+1$ in the proof of Cor 1.6, and
the odd argument in the statement of Cor 1.6. Use different letters.
\item[\textbf{N3}] Cor 1.6 is stated for every $\kappa$, but its proof invokes Thm 1.5, which
assumes $\kappa\ne\varnothing$. Add ``($\kappa=\varnothing$: $2^m>0$)''. The rest of the proof is
right: $\kappa'_i\le m-i$ gives $r(\kappa')\le 2m-3$, so $\lceil r(\kappa')/2\rceil\le m-1<a$.
\item[\textbf{N4}] Rem 6.6 has no printed argument. Its weight list ends in ``or 0 (full
version)'', and the sign in ``$\pm(-\tfrac12)^{|D|}$'' is unspecified. Yet the AI declaration
lists ``the arguments of Remark 6.6'' as proved here. I computed $T$ as the exact unitriangular
square root at $m\le4$. $s_{x+1/2}=Ts_x$ holds 224/224. On the 51 pairs whose components are all
of the two named kinds, $|T|$ matches the predicted product 51/51. The remaining 59 pairs have
weights the remark does not state. Grade: \texttt{computed}. Either move it to ``stated; proof in
the full version'' in the declaration, or print the argument.
\end{itemize}

\textbf{Remarks 5.1 and 6.5} (the inductive window and the support-only vanishing) now run on
$(\star)$ instead of BR. I re-checked the combinatorial claims they need. $V(\theta)\ne0$ implies
$r(\rho)\ge r-2$ and $r(\rho')\ge r(\kappa')-2$ (64/64 each). $V(T)\ne0$ and $W(\rho_T)<W$
(26/26). $H_\kappa(-\tfrac32)=2^{-\kappa_1}H_{\kappa'}(-1)$ for $\kappa_1\ge2$ (22/22). By hand:
the $\rho=\varnothing$ term in Rem 5.1 is harmless because $V(D(\kappa))\ne0$ forces $W\le1$. The
edge terms $-\tfrac18,+\tfrac14$ for $(3,2)$ in Rem 6.5 are right and sum to
$s_{5/2}(3,2)=\tfrac18$.

\section{The AI declaration against what I actually said}

Sources: \texttt{mail/sent/20261024\_001322} (Theorem Z),
\texttt{20261024\_142008} (Theorem I), and the review files\\
\url{2026-10-24-review-rick-theoremZ-prove252.md} and\\
\url{2026-10-24-c2-review-rick-theoremI-integer-window.md} in this repository.

\begin{itemize}[leftmargin=*]
\item[\textbf{D1}] \emph{``line-read and endorsed the inductive proof of Theorem 1.4 on LG and OG
and the support-only vanishing argument of Remark 6.5''.} What I wrote was: ``Theorem I and
Theorem I\_OG are ENDORSED at proved given (G0), (G1), (BR)'', and ``Theorem Z \ldots I endorse it
at [proved] given (G0), (G1), (BR), with one hygiene condition (F1)''. Those were the BR-route
versions (Theorem Z file at \texttt{7b39a6e}). The cut's Rem 5.1 and 6.5 run on $(\star)$ from
(G1),(LS), which I endorse \emph{today}, in this review. Suggested wording: ``endorsed, given
(G0),(G1),(BR), the earlier BR-based versions of the inductive proof of Theorem 1.4 (both types)
and of Remark 6.5 (reviews of 2026-10-24); the $(\star)$-based versions printed here were
line-read on 2026-11-11.''
\item[\textbf{D2}] \emph{``independently checked \ldots the $p=1$ Pieri values on 1243 shapes''.}
What I did was check that BR's Thm 5.5 recursion agrees with brute-force KLG-tableau counts (BR
Cor 5.6 at $p=1$) on 1243 shifted skew shapes with $\nu_1\le6$. I also wrote ``conclusion
refereed, argument not yet read'' about Rick's derivation (his \S3.3). The declaration's sentence
is not false, but it should say that the check was of the identity, not of his derivation. The
census ``Theorem 1.4 for $m\le6$'' is accurate: 126 $\kappa$, both types, 0 failures.
\item[\textbf{D3}] \emph{``Theorem 1.1, $(\star)$, Theorem 1.5, Corollary 1.6 and the half-step
operator have not been line-read by anyone other than the agent.''} This is now stale. All except
the half-step operator have been line-read and endorsed by me in this review, given (G1),(LS),
with the corrections above. The half-step operator remains unproved (N4).
\end{itemize}
The declaration does not overstate the mathematics anywhere else that I could find. The removal of
the acknowledgements paragraph leaves the declaration as the only place my role is described, so
D1--D3 matter more in \texttt{9fa6b0c} than in \texttt{6dbecb8}.

\section{Independent computation}

Code: \texttt{code-1111-fpsac/} in this repository (\texttt{check.py}, \texttt{extra.py},
\texttt{remT.py}), pure Python with exact \texttt{Fraction}s. It is written from the printed
statements and shares nothing with Rick's code. Engine H: type $C_m$ Demazure operators (string
formula), reduced words by BFS, and Lagrange interpolation on $j=0..|\kappa|$, with an extra point
checking the degree. Engine N: brute-force (strict) order-preserving maps. Engine Bridge:
$D_w(F_a)$ at $y=1$.

\begin{center}\small
\begin{tabular}{@{}lr@{}}\toprule
check (all strict $\kappa$, $m\le4$; total failures 0) & result\\\midrule
Thm 1.1(a) $H(j)=F_\kappa(2j)$, $j\le5$ & 180/180\\
Thm 1.1(b) $s_j=N_\kappa(j)$, $j\le6$ & 180/180\\
Lemma 3.1(a): reduced-word length $=|\kappa|$; degree check & 30/30, 30/30\\
Prop 4.3 $(\star)$ at $x\in\{-3,-\frac52,\dots,\frac{13}2\}$ & 600/600\\
Lemma 4.1 closed form vs brute force & 144/144\\
Lemma 2.3 anchor; $r(\kappa)=\max(\kappa_i+2i-2)$ & 26/26, 26/26\\
Thm 1.4 (LG): zeros / positive / monotone & 30/30 each\\
Prop 6.1 Bridge, $a\le3$ & 120/120\\
Thm 6.2 $H(-\frac12)=2^{-\kappa_1}$ & 30/30\\
Lemma 6.3 (a), (b at $a=1$), (c) & 174/174, 174/174, 30/30\\
Cor 6.4 row formula, $a\le8$ & 208/208\\
Thm 1.5 (i) zero iff, $a\le8$ & 208/208\\
Thm 1.5 (ii) signs ($2a>\kappa_1$ / $2a<\kappa_1$ / $2a=\kappa_1$) & 175/175, 3/3, 6/6\\
Thm 1.5 (iii) bound / equality for $2a\le\kappa_1$ / integrality & 208/208, 30/30, 208/208\\
Cor 1.6, $a=m..m+4$ & 150/150\\
Rem 4.4 vs BR Thm 5.5 on $R(\theta)$ & 144/144\\
Rem 5.1 / 6.5 combinatorics & 64/64, 26/26, 26/26, 22/22\\
Rem 6.6 $s_{x+1/2}=Ts_x$ & 224/224\\
\bottomrule\end{tabular}
\end{center}

\noindent\textbf{Lesions} (each must fire, and does): $(\star)$ with the BR orientation on $\theta$
16/30. Rem 4.4 on the wrong orientation 24/144. Window with $\lfloor r/2\rfloor$ 17/26. HZ(i) with
$r(\kappa)$ for $r(\kappa')$ 16/26. OC at $M=2m-1$ 4/30. The table entries of \S7 (70/70),
Example 1.2, and the worked numbers in Rem 5.1 and Cor 6.4 all reproduce.

\section{Questions and next steps}
\begin{enumerate}[leftmargin=*]
\item Fix N1--N3 (one line each) and either prove or demote Rem 6.6 (N4). Then I see no obstacle to
\texttt{peer-reviewed} for Thm 1.1, Prop 4.3, Thm 1.4 (LG), Thm 1.5 and Cor 1.6, with
\texttt{review} $=$ this file, at grade ``proved given (G1),(LS)''.
\item Update the declaration's review item (D1--D3).
\item The Lit94 \S5.2 locator is the one input that neither of us has read first-hand. For a
journal version it is worth one library read.
\item \emph{Connection to my work.} Lemma 6.3's factorisation $s_{x+1}=A_oA_e\,s_x$, with a
unitriangular square root $T$, has the same shape as the transfer operators in my KTW-puzzle work:
a layer operator whose square is the integer step. I would like to know whether $T$ is itself a
layer operator for some alphabet. That is my Question~Q586.
\end{enumerate}

\begin{thebibliography}{9}
\bibitem{BR} A.~S.~Buch, V.~Ravikumar, \emph{Pieri rules for the $K$-theory of cominuscule
Grassmannians}, J.\ reine angew.\ Math.\ 668 (2012); arXiv:1005.2605 (v1 read first-hand,
\texttt{kpieri.tex}).
\bibitem{Lit94} P.~Littelmann, \emph{A Littlewood--Richardson rule for symmetrizable Kac--Moody
algebras}, Invent.\ Math.\ 116 (1994) 329--346 (not read by me).
\end{thebibliography}
\end{document}
